%% main.tex — UnBTeX starter (English mode)
%%
%% ABNT-compliant academic work at the University of Brasilia. ABNT NBR
%% 14724 requires a Portuguese-language abstract (Resumo) even for a work
%% written in English, alongside the English one, so both are present below
%% — this is an institutional requirement, not a mistake.
%%
%% Compile: pdflatex main, bibtex main, pdflatex main, pdflatex main.

\documentclass[
    oneside,
    idioma=english,
    bib=alf,
    refback,
    num=tradicional,
]{unbtex}

\usepackage{lscape}
\usepackage{afterpage}

\graphicspath{{figures}}

\begin{document}

\titulo{Efficient Query Planning for Sensor Time-Series Databases}
\tituloestrangeiro{}

\autori[]{Camila}{Duarte}
\autorii[]{}{}

\numerocutter{}

\orientador[Advisor]{Prof. Dr.}{Ricardo}{Almeida}
\coorientador[Co-advisor]{}{}{}

\tipotrabalho{Master's Dissertation}
\tipocurso{Programa de P\'os-Gradua\c{c}\~ao em Engenharia El\'etrica}
\preambulo{Master's Dissertation submitted to the Graduate Program in Electrical Engineering of the University of Bras\'ilia as part of the requirements for the degree of Master.}
\publicacao{}

\instituicao[Universidade de Bras\'ilia]{Faculdade de Tecnologia}{Departamento de Engenharia El\'etrica}

\local{Bras\'ilia}
\dia{20}
\mes{maio}
\ano{2026}

\membrodabancai{Prof. Dr. Ricardo Almeida\\ UnB/FT/ENE}
\membrodabancaifuncao{Advisor}
\membrodabancaii{Prof. Dr. Fernanda Lima\\ UnB/FT/ENE}
\membrodabancaiifuncao{Internal examiner}
\membrodabancaiii{Prof. Dr. Tiago Novaes\\ UnB/CIC}
\membrodabancaiiifuncao{External examiner}
\membrodabancaiv{}
\membrodabancaivfuncao{}
\membrodabancav{}
\membrodabancavfuncao{}

\begin{Resumo}
Bancos de dados de s\'eries temporais de sensores crescem rapidamente em
volume, o que torna o planejamento eficiente de consultas essencial para
aplica\c{c}\~oes de monitoramento em tempo real. Esta disserta\c{c}\~ao
prop\~oe uma estrat\'egia de planejamento de consultas que explora a
regularidade temporal desses dados para reduzir o custo de varredura,
avaliada sobre tr\^es conjuntos de dados de sensoriamento ambiental de
grande escala.
\end{Resumo}

\begin{Abstract}
Sensor time-series databases grow rapidly in volume, which makes efficient
query planning essential for real-time monitoring applications. This
dissertation proposes a query planning strategy that exploits the temporal
regularity of such data to reduce scan cost, evaluated on three large-scale
environmental sensing datasets. The proposed planner reduces average query
latency substantially relative to a generic range-scan baseline, with the
largest gains on queries that aggregate across many sensors over long time
windows.
\end{Abstract}

\pchavei{Banco de dados}
\pchaveii{S\'eries temporais}
\pchaveiii{Planejamento de consultas}
\pchaveiv{Sensoriamento}
\pchavev{}

\kwordi{Database}
\kwordii{Time series}
\kwordiii{Query planning}
\kwordiv{Sensing}
\kwordv{}

\idiomaagradecimentos{english}

\begin{AgradecimentosAutorI}
I thank my advisor for the guidance and patience shown throughout this
project, and my committee members for their careful reading of this work.
Thanks also to the members of the Distributed Systems Lab for many helpful
discussions, and to the university's High-Performance Computing center for
providing access to the compute resources used in the evaluation.
\end{AgradecimentosAutorI}

\begin{AgradecimentosAutorII}
\end{AgradecimentosAutorII}

\begin{DedicatoriaAutorI}
To everyone who ever waited patiently for a query to finish.
\end{DedicatoriaAutorI}

\begin{DedicatoriaAutorII}
\end{DedicatoriaAutorII}

\pretextual

\imprimirfolhaderosto*

\imprimirfolhadeaprovacao

\imprimirdedicatoria

\imprimiragradecimentos

\imprimirresumos

\imprimirlistadefiguras

\imprimirlistadetabelas

\imprimirsumario

\textual

\chapter{Introduction}
\label{cap:intro}

Sensor deployments for environmental and infrastructure monitoring routinely
produce time-series data at a scale that makes naive range-scan query
execution impractical. A query asking for the average temperature across a
thousand sensors over the last month may need to scan billions of
individual readings if the underlying storage engine has no way to exploit
the temporal and sensor-wise structure of the data \cite{Jensen2017}.

\section{Motivation}
Existing time-series database systems typically rely on generic indexing
structures borrowed from relational systems, which do not exploit the fact
that sensor readings arrive at regular intervals and are naturally
partitioned by sensor identity and time window \cite{Pelkonen2015}.

\section{Contributions}
This dissertation makes the following contributions:
\begin{itemize}
  \item A query planning strategy that partitions scans along both the
    sensor and time dimensions, exploiting known sampling regularity to
    prune irrelevant partitions before execution.
  \item An implementation evaluated against three large-scale environmental
    sensing datasets with heterogeneous sampling rates.
  \item An analysis of when the proposed planner's gains are largest,
    showing the biggest improvements on queries aggregating across many
    sensors over long time windows.
\end{itemize}

\section{Dissertation Outline}
Chapter~\ref{cap:related} reviews prior work on time-series database
systems and query planning. Later chapters (omitted from this starter
document) would present the planning algorithm, implementation, and
experimental results in detail.

\chapter{Related Work}
\label{cap:related}

\section{Time-Series Storage Engines}
Purpose-built time-series databases such as Gorilla store data in
compressed, append-only blocks keyed by sensor and time window, which
improves both write throughput and compression ratio relative to
general-purpose storage engines \cite{Pelkonen2015}.

\section{Query Planning over Temporal Data}
Query optimization for temporal data has a long history in the database
literature, though most classical work targets transactional workloads
rather than the high-cardinality, append-mostly workloads typical of sensor
telemetry \cite{Jensen2017}.

\section{Gap Addressed by This Dissertation}
Prior time-series query planners generally treat sensor identity and time
as independent dimensions to be indexed separately. This dissertation
explores planning jointly over both dimensions when sampling rate is known
in advance.


\postextual

\bibliography{references}

\end{document}
